% Fragmento integrable. Preambulo anfitrion: amsmath, amssymb, longtable, hyperref.
\section{Profondeur nonadique et persistance de l’itinéraire}
\label{hmtfg:fragmento:profundidad-nonadica-y-persistencia-del-itinerario}

\subsection{Croissance extérieure et résolution intérieure}
\label{hmtfg:prolongacion:2}

La loi de raffinement de l’article XI, de base entière \(B\ge2\), est
\[
L_c(x,y)=(Bx-c,By+c),\qquad 0\le c<B.
\]
Pour un mot \(c_1\ldots c_m\), son registre est
\[
C_m=\sum_{j=1}^m c_jB^{m-j},\qquad
(x_m,y_m)=(B^mx_0-C_m,B^my_0+C_m).
\]
La somme extérieure vaut \(B^mM_0\), avec \(M_0=x_0+y_0\). La lecture normalisée du registre appartient au cylindre
\[
I_m(C_m)=\left[\frac{C_m}{B^m},\frac{C_m+1}{B^m}\right],
\qquad \operatorname{diam}I_m=B^{-m}.
\]
Le couple entier et l’échelle restituent les quotients, les restes et les zéros initiaux du mot. Les extrémités à double développement conservent leur marque de frontière. Ainsi se coordonnent la croissance du support et le raffinement de la lecture sans identifier des histoires distinctes par leur seule image limite.

La composition est exacte :
\[
L_D^{[b]}\circ L_C^{[a]}=L_{B^bC+D}^{[a+b]}.
\]
Par associativité, regrouper une histoire de dix-huit événements en neuf blocs de deux ou en deux blocs de neuf produit le même registre. Au niveau de retour \(2^N\), l’horizon commun est \(9\,2^N\). La preuve est une égalité de composition sur la même histoire ordonnée ; elle conserve ses retenues et ses sous-blocs récupérables.

La restriction des états et le prolongement de leurs lecteurs s’expriment par les limites suivantes :
\[
X_\infty=\varprojlim X_m,\qquad
r^*a=a\circ r,\qquad
r^*e_x=\sum_{r(y)=x}e_y,
\qquad \mathfrak H_{\rm cyl}=\varinjlim\mathfrak H_m.
\]
Les états se restreignent ; les lecteurs se prolongent. L’image réciproque préserve les produits, l’unité et les idempotents. Une famille exactement compatible induit un lecteur sur la limite inductive. Les approximants qui ne sont compatibles qu’à une erreur près exigent en outre une estimation de convergence dans leur complétion. La section suivante fournit cette estimation pour le retour critique.

\subsubsection{Capacité et vacances du registre}

Le lecteur de capacité de XI compare des blocs de \(729\) et de \(1000\) possibilités. Avec \(\rho=\log_{1000}729=\log_{10}9\), il conserve
\[
\mathcal C_k=\lfloor(k+1)\rho\rfloor,\qquad
\mathcal C_k-\mathcal C_{k-1}=1-z_k,
\]
\[
\mathcal C_b-\mathcal C_a+\sum_{k=a+1}^b z_k=b-a,
\qquad T_{\rm cap}(a)=\lceil a/\rho\rceil-1\quad(a\ge1).
\]
Chaque vacance enregistre une transition pendant laquelle la capacité se maintient ; la longueur chronologique est conservée dans le bilan. La phase de capacité et la phase chronologique sont des lecteurs distincts.

Pour stocker un préfixe de \(m\) chiffres en base neuf, une capacité décimale suffisante est
\[
a(m)=\left\lceil\frac{m\rho}{3}\right\rceil.
\]
Le quotient et le reste permettent de convertir l’entier \(C_m\) entre bases, en conservant séparément la profondeur \(m\). On a
\[
9^m\le1000^{a(m)}\le729^{T_{\rm cap}(a(m))+1}.
\]
Ces inégalités se décident également au moyen de puissances entières, sans arrondir les logarithmes. Le budget de lecture de la section suivante se traduit ainsi en capacité du registre en conservant les unités. Ce calcul dimensionne le stockage ; les règles de survie déterminent en outre quelles histoires natives sont admissibles.


\subsection{Budget explicite de profondeur pour chaque retour}
\label{hmtfg:prolongacion:3}

\textbf{Théorème.} Pour \(|\eta|\le1/40\) fixé, normalisons
\[
t=\lambda u_\eta(1)/2\in[0,1],\qquad z=(x+1)/2,
\]
\[
F_{t,\eta}(z)=1-t\frac{u_\eta(2z-1)}{u_\eta(1)},\qquad z\in[0,1].
\]
Pour tout entier \(q\ge1\) et deux paramètres \(t,s\),
\[
\left|F_{t,\eta}^{q}(1/2)-F_{s,\eta}^{q}(1/2)\right|
\le S_q|t-s|,
\qquad S_q=\sum_{j=0}^{q-1}(6\sqrt2)^j<9^q.
\]
Par conséquent, un cylindre de paramètres en base neuf, de profondeur \(m=2^N+r\), enferme la lecture du retour \(q=2^N\) dans un intervalle de diamètre inférieur à \(9^{-r}\).

\textbf{Démonstration.} La fonction \((1-\cos v)/v^2\) décroît pour \(0<v<\pi\) : le signe de sa dérivée est celui de \(v\sin v-2(1-\cos v)\), négatif car \(\tan(v/2)>v/2\). Par le deuxième moment,
\[
u_\eta(1)\ge\frac{1-\cos3}{9}\sum_jw_jk_j^2
=\frac{2(1-\cos3)}9>\frac13.
\]
La dernière inégalité se vérifie, par exemple, en majorant \(\cos3\) par sa somme de Taylor alternée jusqu’au degré huit, qui est inférieure à \(-1/2\). De plus, Cauchy–Schwarz donne \(|u_\eta'(x)|\le\sqrt2\) sur la droite réelle. On en déduit
\[
\operatorname{Lip}_zF_{t,\eta}\le6\sqrt2<9,\qquad
\operatorname{Lip}_tF_{t,\eta}\le1.
\]
Si \(d_j\) est la différence entre les deux orbites, alors \(d_0=0\) et \(d_{j+1}\le6\sqrt2\,d_j+|t-s|\). La récurrence prouve la somme géométrique. Pour \(|t-s|\le9^{-m}\), le diamètre est inférieur à \(9^{q-m}\). La substitution de \(q=2^N\) et de \(m=q+r\) achève la preuve. \(\square\)

La borne contrôle l’incertitude du paramètre pour \(\eta\) fixé. Une implémentation ajoute et propage séparément l’erreur d’évaluation des cosinus et l’erreur d’arrondi ; le certificat rationnel antérieur fournit ces opérations. La borne est suffisante et conservatrice, avec un coût qui croît avec l’horizon.

\textbf{Lecteur limite.} Réalisons \(t\) à partir de \((x_0,y_0)=(1,0)\), \(B=9\), de sorte que
\[
(x_m,y_m)=(9^m-C_m,C_m),\qquad t_m=C_m/9^m.
\]
Les lecteurs \(E_{q,m}=F_{t_m,\eta}^{q}(1/2)\), ramenés à un niveau commun, satisfont
\[
\|E_{q,m+a}-r^*E_{q,m}\|_\infty\le S_q9^{-m}.
\]
Ils convergent uniformément vers un lecteur continu \(E_q\) dans la complétion de l’algèbre cylindrique. La version ensembliste
\[
J_{q,m}(C)=\{F_{t,\eta}^{q}(1/2):t\in I_m(C)\}
\]
est exactement emboîtée et ses diamètres tendent vers zéro. Ce résultat donne le passage à la limite de la \textbf{lecture à horizon fixé}, avec contrôle de profondeur à tout horizon demandé.


\subsection{Mot de doublement et lecteur nonadique de la chronologie}
\label{hmtfg:prolongacion:4}

Les huit itinéraires certifiés dans le travail antérieur obéissent à
\[
W_1=+,\qquad W_{n+1}=W_n\,(-1)^n\,W_n.
\]
Son unique prolongement préfixiel est
\[
\tau_j=(-1)^{v_2(j)},\qquad j\ge1;
\qquad \tau_{2j}=-\tau_j,\quad\tau_{2j+1}=+.
\]
La preuve utilise \(v_2(2^n+r)=v_2(r)\) pour \(0<r<2^n\). C’est une récurrence à toute profondeur, et la concordance avec les huit retours calculés constitue sa vérification finie dans la famille.

Dans l’ordre lexicographique à parité d’une application à maximum critique, le facteur d’orientation précédant la position \(k\) est \(\prod_{j<k}(-\tau_j)\). Le mot \(\tau\) est strictement supérieur à chacun de ses décalages positifs. En effet, le premier désaccord avec un décalage impair se produit en \(k=1\) ou \(2\) ; un décalage pair réduit la comparaison de moitié et double \(k\). Le premier désaccord se situe donc en \(k=2^r\). À cet endroit
\[
\prod_{j<2^r}(-\tau_j)=(-1)^r=\tau_{2^r},
\]
et la différence orientée de signe opposé est positive. Cette démonstration prouve également l’apériodicité.

Le registre chronologique conserve
\[
K=\rho_9(K)+9q_9(K).
\]
Son lecteur dyadique est
\[
\ell_n(B_K)=[\rho_9(K)+9q_9(K)]_{2^n}=[K]_{2^n}.
\]
Sur la carte des histoires où \(\Gamma_9 B_K=B_{K+9}\),
\[
\ell_n(\Gamma_9B_K)=\ell_n(B_K)+9\pmod{2^n},\qquad
\pi_{n+1,n}\ell_{n+1}=\ell_n.
\]
Comme neuf est impair, l’avancement de neuf visite les \(2^n\) positions. Sa conjugaison avec l’avancement unitaire est \(j\mapsto9j\). Pour \(0<j<2^n\),
\[
v_2(9j\bmod2^n)=v_2(j),
\]
de sorte que le rééchantillonnage nonadique conserve également le mot orienté dans chaque cycle certifié.

Ce lecteur utilise le reste \textbf{et} le quotient. Les phases de \(K=1\) et de \(K=10\) coïncident modulo neuf, tandis que leurs signes dyadiques sont opposés. La mémoire distingue les deux histoires. La propriété arithmétique du rééchantillonnage vaut pour tout multiplicateur impair ; le calendrier HMT fixe ici le choix de neuf.


\subsection{Mémoire bilatérale compatible à profondeur infinie}
\label{hmtfg:prolongacion:5}

Soient \(\mathcal H_n=\ell^2(\mathbb Z/2^n\mathbb Z)\),
\[
C_ne_j=e_{j+1},\qquad
J_ne_j=(e_j+e_{j+2^n})/\sqrt2.
\]
Une vérification sur la base donne \(J_n^*J_n=I\) et \(C_{n+1}J_n=J_nC_n\), y compris les termes de bord.

L’article XI fournit les projecteurs
\[
P_0=\frac19\begin{pmatrix}8&-\sqrt8\\-\sqrt8&1\end{pmatrix},\qquad
P_1=\frac19\begin{pmatrix}1&\sqrt8\\\sqrt8&8\end{pmatrix}.
\]
Ils sont orthogonaux et complémentaires. Par conséquent
\[
\mathbb U_n=P_0\otimes I+P_1\otimes C_n
=\begin{pmatrix}T_n&D_n\\D_n&R_n\end{pmatrix}
\]
est unitaire, avec
\[
T_n=(8I+C_n)/9,\quad D_n=\sqrt8(C_n-I)/9,\quad R_n=(I+8C_n)/9.
\]
La naturalité déjà démontrée implique
\[
\mathbb U_{n+1}(I_2\otimes J_n)=(I_2\otimes J_n)\mathbb U_n,
\quad \mathbb U_n^a=P_0\otimes I+P_1\otimes C_n^a
\quad(a\in\mathbb Z).
\]
La limite hilbertienne s’identifie à \(L^2(\mathbb Z_2,\mu_{\rm Haar})\), en envoyant \(e_j^{(n)}\) sur \(2^{n/2}\mathbf1_{j+2^n\mathbb Z_2}\). Les identités se prolongent par densité et bornitude. Pour tout vecteur \(v\) et tout temps entier \(a\),
\[
\boxed{\|T_{\infty,a}v\|^2+\|D_{\infty,a}v\|^2=\|v\|^2,}
\]
\[
\boxed{v=T_{\infty,a}^*(T_{\infty,a}v)+D_{\infty,a}^*(D_{\infty,a}v).}
\]
Ici \(T_{\infty,a}=(8I+C_\infty^a)/9\). La composition bilatérale conserve l’archive entre les pas ; composer seulement \(T_n\) décrit une autre opération.

La permutation \(S_ne_j=e_{9j\bmod2^n}\) satisfait
\[
S_{n+1}J_n=J_nS_n,\qquad S_nC_nS_n^{-1}=C_n^9.
\]
On obtient ainsi à la limite
\[
S_\infty C_\infty S_\infty^{-1}=C_\infty^9,
\quad (I_2\otimes S_\infty)\mathbb U_\infty
(I_2\otimes S_\infty^{-1})=\mathbb U_\infty^9.
\]
La représentation temporelle est fidèle : \(C_\infty^a=C_\infty^b\) impose \(2^n\mid a-b\) pour tout \(n\), d’où \(a=b\). Des préparations particulières peuvent conserver des symétries temporelles. L’identité operatorielle conserve cette distinction.

Le point projectif \(([K]_{2^n})_n\) et un vecteur de \(L^2\) sont de types différents : une histoire ponctuelle induit une mesure de Dirac, tandis que les amplitudes normalisables appartiennent à l’espace de Hilbert. Le résultat construit un lecteur équivariant de la chronologie ; les autres coordonnées natives du TPK accompagnent le registre complet.


\subsection{Existence et séparation des obligations asymptotiques}
\label{hmtfg:prolongacion:6}

\subsubsection{Réalisation du mot infini sur toute la bande}
\label{hmtfg:prolongacion:6.1}

\textbf{Théorème d’existence.} Pour chaque \(|\eta|\le1/40\) il existe au moins un
\[
\lambda_\infty(\eta)\in
\left[\frac1{2u_\eta(1)},\frac2{u_\eta(1)}\right]
\]
tel que
\[
\boxed{\operatorname{sign}f_{\lambda_\infty(\eta),\eta}^{j}(0)
=(-1)^{v_2(j)}\quad\text{pour tout }j\ge1.}
\]
La conclusion réalise l’itinéraire infini de doublement au sein de la famille fixée. Le quantificateur d’existence reste séparé d’une sélection unique et de la loi des distances entre paramètres superstables.

\textbf{Démonstration.} Nous fixons \(\eta\). Nous écrivons \(K_\lambda\) pour le mot de signes de l’orbite critique ; lorsque celle-ci revient pour la première fois à zéro, le mot se termine par ce zéro. Nous utilisons l’ordre unimodal à maximum, dont le facteur avant la comparaison du symbole \(j\) est le produit des signes \(-K_i\), \(i<j\).

À l’extrémité gauche, l’image de tout l’intervalle est dans \([1/2,1]\) ; donc \(K_-=+++\cdots\). À l’extrémité droite, l’orbite est \(0\mapsto1\mapsto-1\mapsto-1\), d’où \(K_+=+---\cdots\). La comparaison des deuxième et troisième symboles, respectivement, donne
\[
K_-<\tau<K_+.
\]

Si \(K_{\lambda_0}\) est infini et distinct de \(\tau\), la première différence se produit à une position finie. La continuité des itérées correspondantes rend localement constant le côté de la comparaison. Il reste à examiner un paramètre superstable de première période \(p\), de mot \(W0\), où \(|W|=p-1\).

Pour \(\lambda\) proche de \(\lambda_0\), soient \(g_\lambda=f_{\lambda,\eta}^p\) et \(a=g_\lambda(0)\). On a \(g_\lambda'(0)=0\), et une borne locale uniforme de la dérivée seconde donne
\[
g_\lambda(x)=a+O(x^2).
\]
Pour \(|a|\) suffisamment petit et non nul, \(g_\lambda\) envoie \([-2|a|,2|a|]\) dans \([a-|a|/2,a+|a|/2]\). Les retours conservent le signe de \(a\), et les \(p-1\) positions intermédiaires conservent \(W\). Les mots voisins sont donc \((W+)^\infty\) ou \((W-)^\infty\) ; les paramètres avec \(a=0\) conservent \(W0\).

Si \(\tau\) diffère de \(W\) avant la position \(p\), cette différence fixe le même côté pour les trois mots. Si elle partage \(W\), posons
\[
s=\tau_p,\qquad e=\prod_{j<p}(-\tau_j),\qquad A=Ws.
\]
Le signe de la comparaison de \(\tau\) avec \(W0\) et avec \((W,-s)^\infty\) est \(es\). Pour la comparer à \(A^\infty\), soit \(q\) la première position où \(\tau_{p+q}\ne\tau_q\). Elle existe par apériodicité. Jusqu’à la position \(p+q-1\), les deux mots coïncident, et le signe d’orientation se factorise en \(O_pO_{q-1}\), avec \(O_p=-es\). La maximalité stricte déjà démontrée donne
\[
O_{q-1}(\tau_q-\tau_{p+q})>0.
\]
Par conséquent, la comparaison de \(\tau\) avec \(A^\infty\) a pour signe \(-O_p=es\). Les trois possibilités locales sont du même côté de \(\tau\).

Si aucun paramètre ne réalisait \(\tau\), les ensembles \(\{\lambda:K_\lambda<\tau\}\) et \(\{\lambda:K_\lambda>\tau\}\) seraient ouverts, disjoints, non vides et couvriraient l’intervalle des paramètres. La connexité de l’intervalle exclut cette partition. Le paramètre annoncé existe donc. \(\square\)

Cet argument de continuité des itinéraires s’applique après fixation du générateur HMT et de sa famille analytique ; ses hypothèses sont vérifiées ici directement. La construction conserve le mot infini sans introduire comme donnée un paramètre de la famille quadratique ni la valeur universelle.

\subsubsection{Marges empêchant de perdre l’itinéraire lors du passage à la limite}
\label{hmtfg:prolongacion:6.2}

Écrivons \(c_p=f_{\lambda,\eta}^p(0)\). Une borne rationnelle uniforme suffisante est
\[
|f'|,|f''|\le16,\qquad
B_p:=16^p\frac{16^p-1}{15}\ge\|(f^p)''\|_{[-1,1]}.
\]
En effet, \(1-\cos v\ge v^2/2-v^4/24\ge v^2/8\) pour \(|v|\le3\), de sorte que \(u_\eta(1)\ge1/4\), \(\lambda\le8\) et les deux dérivées sont majorés par \(16\), en utilisant les moments. La règle de dérivation en chaîne prouve par récurrence la borne \(B_p\).

Le mot \(\tau\) satisfait \(\tau_{2p}=-\tau_p\) pour tout \(p\). Pour chacune de ses réalisations, \(c_{2p}\) et \(c_p\) sont de signes opposés. Comme \((f^p)'(0)=0\), Taylor donne
\[
|c_{2p}-c_p|\le\tfrac12B_p|c_p|^2.
\]
Le membre de gauche est supérieur à \(|c_p|\), et donc
\[
\boxed{|c_p|>2/B_p>0.}
\]
Pour chaque horizon fixé, il existe ainsi une marge uniforme qui empêche une suite de réalisations de perdre son signe en convergeant. Dans les coordonnées compactes \((\eta,b)\), \(b=\lambda u_\eta(1)\in[1/2,2]\), l’ensemble des réalisations de \(\tau\) est fermé et compact, et sa projection sur toute la bande de \(\eta\) est surjective.

Le même argument vaut pour un arbre de préfixes : si tous les préfixes de \(\tau\) sont réalisés, une sous-suite dans l’ensemble compact conserve chaque signe, car les préfixes suffisamment longs comprennent simultanément les positions \(p\) et \(2p\). Le théorème d’existence fournit maintenant la non-vacuité nécessaire à cet argument.

\subsubsection{Intervalles invariants de retour à tous les niveaux}
\label{hmtfg:prolongacion:6.3}

\textbf{Théorème des retours emboîtés.} Pour chacune des réalisations du théorème d’existence de la section \ref{hmtfg:prolongacion:6.1}, soient \(c_j=f^j(0)\) et
\[
J_n=\operatorname{conv}\{c_{2^n},c_{2^{n+1}}\},\qquad n\ge1.
\]
Alors \(0\in\operatorname{int}J_n\), \(J_{n+1}\subset J_n\),
\[
f^{2^n}(J_n)=J_n,
\]
et les \(2^n\) images \(J_n,f(J_n),\ldots,f^{2^n-1}(J_n)\) sont disjointes. Les retours normalisés sont unimodaux à maximum, de valeur critique un et de même mot \(\tau\). Ces intervalles sont des noyaux invariants de retour ; une normalisation exigeant des extrémités périodiques peut élargir le noyau à l’intervalle restrictif maximal correspondant et doit préciser cette étape.

\textbf{Démonstration du premier retour.} Soit \(F\) une application unimodale à maximum, avec \(F(0)=1\) et de mot \(\tau\), et \(d_j=F^j(0)\). Ses premiers signes sont \(+,-,+,+,+,-,+\). L’intervalle \(I=[d_2,1]\) est invariant : \(F(d_2)=d_3>0\), \(F(1)=d_2<0\), et le maximum vaut un. Posons \(J=[d_2,d_4]\). Comme \(F\) décroît sur les arguments positifs et
\[
F(d_3)=d_4>0> d_6=F(d_5),
\]
on a \(d_3<d_5\), et donc \(F(J)=[d_3,1]\).

De plus, \(d_3>d_4\). Pour le vérifier, \(F^2\) est strictement croissante entre les deux points positifs \(d_3,d_4\), car leurs images \(d_4,d_5\) sont positives. Ses valeurs sont \(F^2(d_3)=d_5>0>d_6=F^2(d_4)\), ce qui impose cet ordre. Ainsi \(J\) et \(F(J)\) sont séparés par un intervalle ouvert, et
\[
F^2(J)=F([d_3,1])=[d_2,d_4]=J.
\]
L’application \(F^2|_J\) possède un unique minimum critique : l’image \(F(J)\) reste positive. La conjugaison \(h(x)=d_2x\) inverse l’orientation et normalise le retour en un maximum de valeur un. Son itinéraire est
\[
\operatorname{sign}\frac{d_{2j}}{d_2}
=-\tau_{2j}=\tau_j.
\]
Le même argument peut être répété à toute profondeur.

\textbf{Disjonction par récurrence.} Supposons disjointes les \(p=2^n\) images du parent \(J_n\), et soit \(g=f^p\). Le premier retour appliqué à l’application normalisée produit \(B=J_{n+1}\), \(A=g(B)\), avec \(A\cap B=\varnothing\), \(g(B)=A\), \(g(A)=B\). Pour \(0\le j<p\), une intersection de \(f^j(A)\) et de \(f^j(B)\) produirait, par application de \(f^{p-j}\), une intersection de \(B\) et de \(A\). Ces deux images sont donc disjointes. Les images correspondant à des \(j\) distincts appartiennent à des images disjointes du parent. Les \(2p\) images disjointes et le retour \(f^{2p}(B)=B\) sont ainsi démontrés. La formule des extrémités résulte de la composition des normalisations \(x\mapsto c_{2^n}x\). \(\square\)

Les ensembles compacts
\[
\Lambda_n=\bigcup_{0\le j<2^n}f^j(J_n)
\]
sont non vides et emboîtés. Leur intersection \(\Lambda_\infty\) possède un lecteur continu sur l’odomètre dyadique : à chaque point est attribuée l’étiquette de l’unique composante qu’il occupe à chaque niveau. Toute suite compatible d’étiquettes possède une préimage par compacité, et l’avancement par \(f\) ajoute un aux étiquettes. Cette application est une semi-conjugaison surjective. Sa promotion en conjugaison ponctuelle exige de démontrer que les diamètres des composantes tendent vers zéro ; cette promotion métrique reste séparée.
